% SPDX-License-Identifier: Apache-2.0
\section{The build}
\label{sec:build}

Bazel runs everything in this report, and it fetches every tool it
uses by checksum or by commit. Nothing is installed on the machine
beyond Bazelisk and Git. Table~\ref{tab:tools} lists the tools.

\begin{table}[t]
\centering
\caption{What the build fetches}
\label{tab:tools}
\footnotesize
\begin{tabular}{@{}lp{0.6\columnwidth}@{}}
\toprule
Tool & Version and source \\
\midrule
Vreteno & TxHDL commit \code{\TxhdlCommit}, from its GitHub mirror \\
Tests & riscv-arch-test commit \code{\ActCommit} \\
Reference & \SailVersion, prebuilt release \\
Compiler & \GccVersion \\
Simulator & Verilator 5.046, TxHDL's prebuilt binary \\
C and C++ & LLVM over a pinned Debian sysroot, from TxHDL \\
\LaTeX & a pinned TinyTeX, with TxHDL's copies of IEEEtran, pgf and
listings \\
\bottomrule
\end{tabular}
\end{table}

\subsection{From test file to target}

A repository rule downloads the test suite and reads the configuration
header of every test file. It writes the headers into a list that the
build reads when it analyses the targets, so that selecting the tests
needs no Python. For each selected test, a Bazel macro makes one
target that builds the self-checking ELF in the four steps of
Fig.~\ref{fig:flow}. It makes two more targets that run the ELF, one on
Sail and one on the netlist, and a test target for each run.
\code{sigfmt}, the tool in step three, is a short Rust program. It
replaces the Python function in ACT4 that does the same job.

Each step is an action with declared inputs and outputs, so Bazel runs
the \NTests{} tests in parallel and caches every step. A change to the
core reruns the netlist runs and nothing else. A change to the
description of the core rebuilds every ELF.

\subsection{Known failures}

A test that the core is known to fail is listed with the URL of the
issue that tracks it. Its test target then passes while the run fails,
and fails once the run passes. The list therefore shrinks when the
core is fixed, and \code{bazel test} stays useful for changes to the
harness while a defect is open. This report reads the result files,
not the test targets, so a known failure counts as a failure here.

\subsection{Every day}

A workflow runs the build every day against the newest commit of
TxHDL. It typesets this report, publishes it as a release of the
repository, and publishes the same results as a web page on
hdlfactory.com. The commit of the core and the date are in the title
footnote of every copy.
